\documentclass[10pt,letterpaper]{article}
\usepackage{cornell-notes}

\title{Clause 30.03.01.02.03: Class Locale Id}
\author{}
\date{}
\setDocTitle{Clause 30.03.01.02.03: Class Locale Id}
\setDocSubtitle{C++ 2024}
\setDocOwner{C++ 2024 Cornell Notes}
\setDocDate{\today}
\setCornellCollection{C++ 2024 Cornell Notes}
\setCornellUnitType{Clause}
\setCornellUnitNumber{30.03.01.02.03}
\setCornellUnitTitle{Class Locale Id}
\hypersetup{pdftitle={Clause 30.03.01.02.03: Class Locale Id},pdfsubject={Cornell notes},pdfauthor={}}

\begin{document}
\maketitle
\begin{CornellOverviewBox}{Overview}
\textbf{Classification:} Standard library - Normative\\[2pt]
\textbf{Learning objective:} Understand the rules, terminology, and semantic consequences associated with Class locale::id.
\end{CornellOverviewBox}\begin{CornellSummaryBox}{Summary}
{\sffamily\bfseries\color{Accent} Summary}\par\vspace{3pt}Section 30.3.1.2.3 focuses on Class locale::id. Use the listed grammar productions together with the semantic constraints and disambiguation rules referenced by the section. When applying it, preserve the distinction between mandatory requirements, recommendations, permissions, and behavior deliberately left undefined, unspecified, or implementation-defined.
\end{CornellSummaryBox}

\end{document}
